% Maquette 2 : enregistrement d'un appareil (réparateur, smartphone)
\begin{tikzpicture}
  \draw[mq ecran] (0,0) rectangle (6.6,-11.6);
  \fill[mq barre] (0.08,-0.08) [rounded corners=7pt] -- (0.08,-0.95) [sharp corners] -- (6.52,-0.95) [rounded corners=7pt] -- (6.52,-0.08) [sharp corners] -- cycle;
  \node[mq titre] at (0.35,-0.33) {Qribna};
  \node[mq texte, anchor=north east] at (6.25,-0.33) {Boutique A};

  \node[mq titre] at (0.35,-1.25) {Nouvel appareil};
  \node[mq champ, minimum width=5.9cm] at (0.35,-1.75) {Nom du client};
  \node[mq champ, minimum width=5.9cm] at (0.35,-2.5) {Numéro WhatsApp};
  \node[mq champ, minimum width=5.9cm] at (0.35,-3.25) {Modèle du téléphone};
  \node[mq champ, minimum width=5.9cm, minimum height=1.3cm, anchor=north west,
        text depth=0.75cm] at (0.35,-4.0) {Description de la panne};

  \node[mq titre] at (0.35,-5.65) {Photos de l’état initial};
  \mqphoto{0.35}{-6.1}{1.75}
  \mqphoto{2.43}{-6.1}{1.75}
  \draw[mq cadre, dashed] (4.5,-6.1) rectangle ++(1.75,-1.75);
  \node[mq texte, anchor=center, align=center] at (5.375,-6.975) {+\\Ajouter};

  \node[mq gris] at (0.35,-8.25) {Le numéro de suivi est attribué\\après validation par le client.};

  \node[mq bouton, minimum width=5.9cm] at (3.3,-10.85) {Envoyer pour validation};
\end{tikzpicture}
